%==============================================================
%  lecons/algebre/sous_groupes_additifs.tex — Sous-groupes additifs de R
%==============================================================

\lecontitle{Sous-groupes additifs de \texorpdfstring{$\mathbb{R}$}{R}}{Algèbre \& topologie — Dichotomie monogène ou dense}

%=============================================================
%  § 1 — ÉNONCÉ ET SIGNIFICATION
%=============================================================
\secnum{navyblue}{1}{Énoncé et signification}

\begin{tcolorbox}[thmbox]
  \textbf{Théorème (dichotomie).}\enspace%
  \index{sous-groupe additif de $\mathbb{R}$}\index{sous-groupe!monogène}\index{sous-groupe!dense}
  Soit $G$ un sous-groupe de $(\mathbb{R},+)$. Alors l'une exactement des deux
  situations suivantes se produit :
  \begin{enumerate}
    \item \textbf{cas monogène :} il existe un unique réel $a \geq 0$ tel que
          \[
            G = a\mathbb{Z} = \{\,na : n \in \mathbb{Z}\,\} ;
          \]
    \item \textbf{cas dense :} $G$ est dense dans $\mathbb{R}$.
  \end{enumerate}
  Plus précisément, lorsque $G \neq \{0\}$ on pose
  $a = \inf\bigl(G \cap \,]0,+\infty[\,\bigr)$ : on est dans le cas monogène
  lorsque $a > 0$, et dans le cas dense lorsque $a = 0$. Le cas $G = \{0\}$ est
  monogène de générateur $a = 0$ (la convention $\inf(\emptyset) = +\infty$ sert
  seulement à exprimer qu'alors $G$ n'a aucun élément $> 0$).
\end{tcolorbox}

\begin{significationintuitive}
Un sous-groupe additif de $\mathbb{R}$ ne peut pas être « à moitié discret ».
Soit ses éléments strictement positifs admettent un plus petit représentant $a$,
et alors tout le groupe est constitué des multiples entiers de ce « pas »~; soit
ils s'accumulent en $0$, et la translation par ces éléments arbitrairement petits
permet d'atteindre n'importe quel réel d'aussi près que l'on veut. La frontière
entre les deux mondes est nette : il n'existe pas d'intermédiaire.
\end{significationintuitive}

\decsep{}

%=============================================================
%  § 2 — DÉMONSTRATION
%=============================================================
\secnum{navyblue}{2}{Démonstration}

\begin{ideepreuve}
On mesure la « finesse » du groupe par $a = \inf(G \cap \,]0,+\infty[)$.
Si $a > 0$, on montre d'abord que cette borne inférieure est atteinte
($a \in G$) à l'aide d'un argument de différence de deux éléments trop proches,
puis la division euclidienne réelle donne $G = a\mathbb{Z}$. Si $a = 0$, le
groupe contient des éléments positifs aussi petits que voulu, et leurs multiples
quadrillent $\mathbb{R}$ : c'est la densité.
\end{ideepreuve}

\rubriq{Preuve}

\begin{mdframed}[style=proofbar]
  \textit{Le cas trivial.}\enspace
  Si $G = \{0\}$, alors $G = 0\cdot\mathbb{Z}$ : c'est le cas monogène avec $a=0$.
  On suppose désormais $G \neq \{0\}$. Si $x \in G \setminus \{0\}$, alors $-x \in G$
  (sous-groupe), donc $|x| \in G \cap \,]0,+\infty[$ : cet ensemble est non vide
  et minoré par $0$, sa borne inférieure
  \[
    a = \inf\bigl(G \cap \,]0,+\infty[\,\bigr) \geq 0
  \]
  est bien définie.

  \medskip
  \textit{Cas $a > 0$ : montrons d'abord $a \in G$.}\enspace
  Supposons par l'absurde $a \notin G$. Par définition de la borne inférieure
  appliquée à $\varepsilon = a$, il existe $g \in G$ avec $a < g < 2a$ (l'inégalité
  est stricte à gauche car $a \notin G$). Appliquée ensuite à $\varepsilon = g - a > 0$,
  elle fournit $h \in G$ avec $a < h < g$. Alors
  \[
    g - h \in G \quad\text{et}\quad 0 < g - h < g - a < a,
  \]
  donc $g - h \in G \cap \,]0,+\infty[$ tout en vérifiant $g - h < a$, ce qui
  contredit le fait que $a$ minore $G \cap \,]0,+\infty[$. Donc $a \in G$.

  \medskip
  \textit{Cas $a > 0$ : conclusion $G = a\mathbb{Z}$.}\enspace
  Comme $a \in G$, on a $a\mathbb{Z} \subseteq G$. Réciproquement, soit $x \in G$.
  La division euclidienne réelle donne $q = \lfloor x/a \rfloor \in \mathbb{Z}$ et
  un reste $r = x - qa$ avec $0 \leq r < a$. Or $r = x - qa \in G$ (car $x \in G$
  et $qa \in G$), et $a$ est le plus petit élément strictement positif de $G$ ;
  l'encadrement $0 \leq r < a$ force $r = 0$. Ainsi $x = qa \in a\mathbb{Z}$, d'où
  $G = a\mathbb{Z}$.

  \medskip
  \textit{Cas $a = 0$ : densité.}\enspace
  Soit $x \in \mathbb{R}$ et $\varepsilon > 0$. Puisque $\inf(G \cap \,]0,+\infty[) = 0$,
  il existe $g \in G$ avec $0 < g < \varepsilon$. Posons $n = \lfloor x/g \rfloor \in \mathbb{Z}$,
  de sorte que $ng \leq x < (n+1)g$, donc
  \[
    |x - ng| < g < \varepsilon, \qquad ng \in G.
  \]
  Tout réel est donc limite d'éléments de $G$ : $G$ est dense dans $\mathbb{R}$.

  \medskip
  \textit{Exclusivité et unicité.}\enspace
  Si $a > 0$, l'intervalle $]0,a[$ ne rencontre pas $G$, donc $G$ n'est pas dense :
  les deux cas s'excluent. Enfin, si $a\mathbb{Z} = b\mathbb{Z}$ avec $a,b \geq 0$ :
  si $a = 0$, alors $b\mathbb{Z} = \{0\}$, donc $b = 0$ ; si $a > 0$, alors $a$ est
  le plus petit élément strictement positif de $a\mathbb{Z} = b\mathbb{Z}$, et $b$
  l'est aussi, d'où $a = b$ :
  le générateur positif est unique. \hfill$\blacksquare$
\end{mdframed}

\decsep{}

%=============================================================
%  § 3 — APPLICATION : LE GROUPE $\mathbb{Z} + \alpha\mathbb{Z}$
%=============================================================
\secnum{navyblue}{3}{Application : le groupe \texorpdfstring{$\mathbb{Z}+\alpha\mathbb{Z}$}{Z + aZ}}

\begin{tcolorbox}[thmbox]
  \textbf{Théorème (critère d'irrationalité).}\enspace%
  \index{groupe!$\mathbb{Z}+\alpha\mathbb{Z}$}\index{Kronecker (théorème de)}
  Soit $\alpha \in \mathbb{R}$ et $G = \mathbb{Z} + \alpha\mathbb{Z}
  = \{\,m + n\alpha : (m,n) \in \mathbb{Z}^2\,\}$. Alors :
  \[
    G \text{ est dense dans } \mathbb{R} \iff \alpha \notin \mathbb{Q}.
  \]
  De façon équivalente, la suite des parties fractionnaires $(\{n\alpha\})_{n\in\mathbb{N}}$
  est dense dans $[0,1]$ si et seulement si $\alpha$ est irrationnel.
\end{tcolorbox}

\rubriq{Preuve}

\begin{mdframed}[style=proofbar]
  $G$ est un sous-groupe de $(\mathbb{R},+)$, on lui applique la dichotomie.

  \smallskip\noindent
  \textit{Si $\alpha = p/q \in \mathbb{Q}$} (fraction irréductible, $q \geq 1$),
  alors $m + n\alpha = (mq + np)/q \in \tfrac{1}{q}\mathbb{Z}$. Réciproquement,
  comme $\pgcd(p,q) = 1$, le théorème de Bézout fournit $u,v \in \mathbb{Z}$ avec
  $up + vq = 1$, d'où $\tfrac1q = u\alpha + v \in G$ et donc
  $\tfrac1q\mathbb{Z} \subseteq G$ ; ainsi $G = \tfrac1q\mathbb{Z}$ est monogène,
  \emph{non} dense.

  \smallskip\noindent
  \textit{Si $\alpha \notin \mathbb{Q}$}, supposons $G$ monogène, $G = a\mathbb{Z}$
  avec $a > 0$. Comme $1 \in G$ et $\alpha \in G$, il existe $k, \ell \in \mathbb{Z}$
  avec $1 = ka$ et $\alpha = \ell a$. Alors $a = 1/k$ puis $\alpha = \ell/k \in \mathbb{Q}$,
  contradiction. Le cas monogène est donc exclu : $G$ est dense.

  \smallskip\noindent
  Enfin, chaque $\{n\alpha\} = n\alpha - \lfloor n\alpha\rfloor$ appartient à
  $G \cap [0,1[$ (car $n\alpha \in G$ et $\lfloor n\alpha\rfloor \in \mathbb{Z}
  \subseteq G$). Si $\alpha = p/q \in \mathbb{Q}$, cet ensemble
  $G \cap [0,1[\, = \tfrac1q\mathbb{Z} \cap [0,1[$ est fini : la suite n'est pas
  dense. Si $\alpha \notin \mathbb{Q}$, soit $\varepsilon \in \,]0,1[$ ; par
  densité de $G$, il existe $g = m + r\alpha \in G$ avec $0 < g < \varepsilon$,
  et $r \neq 0$ (sinon $g = m$ serait un entier de $]0,1[$). Si $r > 0$, alors
  $\{r\alpha\} = g$ et, pour tout $k \in \mathbb{N}$ tel que $kg < 1$, on a
  $\{kr\alpha\} = kg$ : ces points quadrillent $[0,1]$ avec un pas $g < \varepsilon$.
  Si $r < 0$, alors $\{(-r)\alpha\} = 1 - g$ et $\{k(-r)\alpha\} = 1 - kg$ tant
  que $kg < 1$ : même quadrillage, parcouru en sens inverse. Dans les deux cas,
  tout point de $[0,1]$ est à distance $< \varepsilon$ d'un $\{n\alpha\}$ avec
  $n \in \mathbb{N}$ : la suite $(\{n\alpha\})_{n\in\mathbb{N}}$ est dense dans
  $[0,1]$.
  \hfill$\blacksquare$
\end{mdframed}

\rubriq{Exemples}

\begin{mdframed}[style=exbar]
  \textbf{1. Pas régulier.}\enspace%
  $G = \frac{1}{2}\mathbb{Z} = \{\ldots,-1,-\tfrac12,0,\tfrac12,1,\ldots\}$ est
  monogène de générateur $a = \tfrac12$.

  \smallskip\noindent
  \textbf{2. $\mathbb{Q}$ est dense.}\enspace%
  $\mathbb{Q}$ est un sous-groupe de $(\mathbb{R},+)$ contenant des éléments
  positifs arbitrairement petits ($\tfrac1n \to 0$) : $a = 0$, il est dense
  --- ce que l'on savait déjà, ici relu par la dichotomie.

  \smallskip\noindent
  \textbf{3. Densité de $(\cos n)_{n\in\mathbb{N}}$.}\enspace%
  \index{densité!de $(\cos n)$}
  Comme $\pi \notin \mathbb{Q}$ (irrationalité admise), $2\pi$ est irrationnel, donc
  $\mathbb{Z} + 2\pi\mathbb{Z}$ est dense dans $\mathbb{R}$ : les entiers $n$ pris
  modulo $2\pi$ sont denses dans $[0,2\pi[$. Par continuité du cosinus,
  $\{\cos n : n \in \mathbb{N}\}$ est dense dans $[-1,1]$.
\end{mdframed}

\rubriq{Contre-exemples et points de vigilance}

\begin{mdframed}[style=cexbar]
  \textbf{$\mathbb{N}$ n'est pas un sous-groupe.}\enspace%
  $\mathbb{N}$ est discret mais n'est \emph{pas} de la forme $a\mathbb{Z}$ : la
  dichotomie ne s'y applique pas, faute d'opposés. L'hypothèse « sous-groupe »
  est essentielle.

  \smallskip\noindent
  \textbf{Dense $\neq$ égal à $\mathbb{R}$.}\enspace%
  $\mathbb{Q}$ et $\mathbb{Z}+\sqrt2\,\mathbb{Z}$ sont denses dans $\mathbb{R}$
  sans être $\mathbb{R}$ tout entier : la densité est une propriété topologique,
  pas une égalité ensembliste.

  \smallskip\noindent
  \textbf{Pas d'intermédiaire.}\enspace%
  Il n'existe pas de sous-groupe de $\mathbb{R}$ qui soit à la fois non monogène
  et non dense : un sous-groupe additif de $\mathbb{R}$ est soit fermé (et alors
  $\{0\}$, $a\mathbb{Z}$ ou $\mathbb{R}$), soit dense.
\end{mdframed}

\exolabel

\begin{mdframed}[style=exobar]
  \begin{enumerate}
    \item Montrer que tout sous-groupe \emph{fermé} de $(\mathbb{R},+)$ est
          $\{0\}$, $a\mathbb{Z}$ ($a>0$), ou $\mathbb{R}$.
    \item Soit $\alpha$ irrationnel. Montrer que $\{m + n\alpha : m,n \in \mathbb{Z}\}$
          rencontre tout intervalle ouvert non vide.
    \item En déduire que $(\sin n)_{n \in \mathbb{N}}$ est dense dans $[-1,1]$.
    \item Montrer que l'adhérence d'un sous-groupe de $(\mathbb{R},+)$ est encore
          un sous-groupe, puis retrouver la dichotomie à partir de ce fait.
    \item Soient $a, b > 0$. Montrer que $a\mathbb{Z} + b\mathbb{Z}$ est monogène
          si et seulement si $a/b \in \mathbb{Q}$, et préciser alors son générateur
          ($\pgcd$ généralisé).
  \end{enumerate}
\end{mdframed}

\leconfooter{Dichotomie des sous-groupes de $(\mathbb{R},+)$ — L.\ Kronecker (1884, densité de $\mathbb{Z}+\alpha\mathbb{Z}$)}
